\documentclass[11pt]{article}
\usepackage[T1]{fontenc}
\usepackage{lmodern,amsmath,amssymb,amsthm,microtype}
\usepackage[margin=1in]{geometry}
\usepackage[colorlinks=true,urlcolor=blue,citecolor=blue,linkcolor=blue]{hyperref}
\newtheorem{theorem}{Theorem}
\newtheorem{corollary}{Corollary}
\newtheorem{remark}{Remark}
\newcommand{\cat}{\mathbin{\Vert}}
\title{Divisibility of Reversed Consecutive-Integer Concatenations\\and an Arbitrary-Permutation Extension}
\author{Qichao Wang\\\small Hebei University of Technology, Tianjin, China}
\date{October 16, 2026}
\begin{document}
\maketitle
\begin{abstract}
For every positive integer $n$, we show that the concatenation of some
$n$ consecutive positive integers in decreasing order is divisible by $n$.
Consequently, OEIS A083468 is positive and OEIS A083469 is well-defined
for every positive index.  We prove a more general assertion: in any even
integer base, \emph{every prescribed permutation} of $n$ consecutive
integers admits infinitely many equal-width concatenations divisible by
$n$.  The proof is constructive and rests only on congruences and the
Chinese remainder theorem.  We establish existence, not an explicit
formula for the least sequence values.
\end{abstract}
\section{The two sequences}
Write $u\cat_b v$ for concatenation of the base-$b$ digit strings
of positive integers $u,v$.  The sequence
\href{https://oeis.org/A083468}{OEIS A083468} gives the smallest
positive multiple of $n$ obtained by concatenating $n$ consecutive
positive integers in \emph{decreasing} order, with the value zero
if no such concatenation exists.  The companion sequence
\href{https://oeis.org/A083469}{OEIS A083469} divides that
smallest value by $n$; its definition does not assign a value in a
hypothetical nonexistence case \cite{oeis468,oeis469}.
We prove nonexistence never occurs.  The proof is independent of the
minimality of the listed sequence terms.  It can also be viewed as a
corollary of the author's even-base arithmetic-progression divisibility
theorem, taking step $-1$ \cite{wangAP}.  The independent proof below
exposes a permutation strengthening.
\section{Divisibility in every prescribed block order}
\begin{theorem}\label{thm:permutation}
Let $b\ge2$ be an even integer, $n\ge1$, and let
$\pi:\{1,\ldots,n\}\to\{0,\ldots,n-1\}$ be an arbitrary bijection.
There exist infinitely many pairs $(d,k)$ of positive integers such that
$k,k+1,\ldots,k+n-1$ all have exactly $d$ digits in base $b$ and
\begin{equation}\label{eq:target}
 n\ \Big|\ (k+\pi(1))\cat_b(k+\pi(2))\cat_b\cdots\cat_b(k+\pi(n)).
\end{equation}
In particular, there are infinitely many distinct concatenated integers
satisfying \eqref{eq:target}.
\end{theorem}
\begin{proof}
Factor $n=qm$, where $q$ is the largest divisor of $n$ supported
on primes dividing $b$. Thus $\gcd(q,m)=1$ and $\gcd(m,b)=1$.
Since $b$ is even, $m$ is odd.
Choose an arbitrarily large positive integer $d$ with
\begin{equation}\label{eq:base}
 q\mid b^d,\qquad b^d\equiv1\pmod m,
 \qquad (b-1)b^{d-1}\ge q+n-1.
\end{equation}
Such $d$ exist arbitrarily far out: divisibility by $q$ holds for all
sufficiently large $d$; the second condition holds when $d$ is a
multiple of the multiplicative order of $b$ modulo $m$ (and is
vacuous when $m=1$); the inequality follows for large $d$.
The inclusive interval
\[
 I_d=[b^{d-1},\,b^d-n]\cap\mathbb Z
\]
contains at least $q$ consecutive integers, by \eqref{eq:base}.
Choose $k\in I_d$ satisfying
\begin{equation}\label{eq:k}
 k+\pi(n)\equiv0\pmod q.
\end{equation}
If $q=1$, this congruence is vacuous. All $n$ of the integers
$k,\ldots,k+n-1$ are positive and have exactly $d$ base-$b$ digits.
Put $B=b^d$. The proposed concatenation is the integer
\begin{equation}\label{eq:poly}
 C=\sum_{i=1}^{n}\bigl(k+\pi(i)\bigr)B^{n-i}.
\end{equation}
Since $q\mid B$, the final summand in \eqref{eq:poly} is the only
one that survives modulo $q$, so \eqref{eq:k} gives
\[
 C\equiv k+\pi(n)\equiv0\pmod q.
\]
On the other hand, $B\equiv1\pmod m$, hence the permutation drops
out of the residue:
\[
 C\equiv\sum_{i=1}^{n}\bigl(k+\pi(i)\bigr)
  =nk+\frac{n(n-1)}2\equiv0\pmod m.
\]
The last congruence holds because $m\mid n$ and $m$ is odd.
As $\gcd(q,m)=1$, the Chinese remainder theorem implies $n\mid C$.
There are arbitrarily large admissible $d$, and the resulting
concatenations have increasing lengths $nd$, proving infinitude.
\end{proof}
\begin{corollary}\label{cor:oeis}
For every positive integer $n$, there exist infinitely many
concatenations of $n$ consecutive positive integers in decreasing
order that are divisible by $n$. Consequently
\[
 \boxed{\mathrm{A083468}(n)>0\quad\text{for every }n\ge1,}
\]
while $\mathrm{A083469}(n)=\mathrm{A083468}(n)/n$ is
well-defined and positive for every $n\ge1$.
\end{corollary}
\begin{proof}
In Theorem~\ref{thm:permutation} take $b=10$ and
$\pi(i)=n-i$, so that the blocks occur in descending order.
\end{proof}
\section{A concrete witness construction}
For the decimal case, factor
\[
 n=q m,\qquad q=2^{v_2(n)}5^{v_5(n)},\qquad \gcd(m,10)=1.
\]
Take an arbitrarily large $d$ satisfying
\[
 q\mid10^d,\qquad10^d\equiv1\pmod m,\qquad
 9\cdot10^{d-1}\ge q+n-1.
\]
The smallest multiple of $q$ not below $10^{d-1}$,
\[
 a=q\left\lceil\frac{10^{d-1}}q\right\rceil,
\]
lies in $[10^{d-1},10^d-n]$. Therefore the integer
\[
 (a+n-1)\cat_{10}(a+n-2)\cat_{10}\cdots\cat_{10}a
\]
is divisible by $n$. For instance, with $n=4$, $d=2$ and $a=12$,
we obtain $15141312$, which is divisible by $4$; this is not the
minimum, which is $5432$ \cite{oeis468}.
\begin{remark}
The hypothesis that $b$ is even ensures that $m$ is odd, which is
used to annihilate $n(n-1)/2$ modulo $m$. No such theorem for arbitrary
odd bases is claimed here.
\end{remark}
\section{Computational checks and scope}
The accompanying script \texttt{verify\_reverse\_concatenation.py}
constructs witnesses for $1\le n\le512$, verifies divisibility using
modular concatenation, and independently reproduces the first sixteen
A083468 minima. Finite checks are supplementary: the theorem is an
unconditional proof for all indices.
The proof establishes \emph{existence} and provides an explicit family of
witnesses. It does not evaluate the minimal concatenation or its
quotient in closed form. The theorem includes the previously observed
reverse-progression consequence of \cite{wangAP}; no claim is made
that the arbitrary-permutation strengthening is new to the literature.
\begin{thebibliography}{9}
\bibitem{oeis468} OEIS Foundation Inc.,
  \emph{A083468}, \url{https://oeis.org/A083468}
  (accessed October 16, 2026).
\bibitem{oeis469} OEIS Foundation Inc.,
  \emph{A083469}, \url{https://oeis.org/A083469}
  (accessed October 16, 2026).
\bibitem{wangAP} Q. Wang,
  \emph{Divisibility of Concatenated Arithmetic Progressions and
  Champernowne Prefixes},
  \url{https://github.com/dongxuelian2/oeis-papers/tree/main/papers/A075000-A075002},
  2026. See the even-base AP divisibility theorem.
\bibitem{luca} F. Luca and B. Teguia Tabuguia,
 \emph{Concatenations of terms of an arithmetic progression},
 Journal of Difference Equations and Applications \textbf{32} (2026), 437--456,
 \url{https://doi.org/10.1080/10236198.2025.2603385}.
\end{thebibliography}
\end{document}
